% Standard preamble for literate programming documents
% Copy this file verbatim to doc/preamble.tex in new projects
%
% This preamble is shared across all literate programming projects.
% The main document (e.g., doc/main.tex) should % Standard preamble for literate programming documents
% Copy this file verbatim to doc/preamble.tex in new projects
%
% This preamble is shared across all literate programming projects.
% The main document (e.g., doc/main.tex) should % Standard preamble for literate programming documents
% Copy this file verbatim to doc/preamble.tex in new projects
%
% This preamble is shared across all literate programming projects.
% The main document (e.g., doc/main.tex) should % Standard preamble for literate programming documents
% Copy this file verbatim to doc/preamble.tex in new projects
%
% This preamble is shared across all literate programming projects.
% The main document (e.g., doc/main.tex) should \input{preamble} after
% \documentclass but before \begin{document}.

%\usepackage[utf8]{inputenc}
\usepackage[T1]{fontenc}
\usepackage[swedish,british]{babel}
\usepackage{booktabs}

\usepackage{noweb}
% Caveat about longxref: it typesets "This definition is continued in
% chunks ..." lists whose line count depends on the page numbers of later
% chunks.  With many continuations (e.g. <<test functions>>) the list
% length can feed back into the pagination it describes, so the labels
% never converge and latexmk reruns forever ("Label(s) may have
% changed.").  If that happens, adjust the prose near the oscillating
% page boundary (diff two consecutive .aux files to find it) until the
% layout has a fixed point again.
\noweboptions{shift,breakcode,longxref,longchunks}

\usepackage[natbib,style=alphabetic,maxbibnames=99]{biblatex}
\addbibresource{bibliography.bib}

\usepackage[inline]{enumitem}
\setlist[enumerate]{label=(\arabic*)}

\usepackage[all]{foreign}
\renewcommand{\foreignfullfont}{}
\renewcommand{\foreignabbrfont}{}

%\usepackage{newclude}
\usepackage{import}

\usepackage[strict]{csquotes}
\usepackage[single]{acro}

\usepackage{subcaption}
\usepackage{graphicx}

\usepackage[noend]{algpseudocode}
\usepackage{xparse}

\let\email\texttt

%\usepackage[outputdir=ltxobj]{minted}
\usepackage{minted}
\setminted{autogobble,linenos}

\usepackage{pythontex}
\setpythontexoutputdir{.}
\setpythontexworkingdir{..}

\usepackage{fancyvrb}

\usepackage{amsmath}
\usepackage{amssymb}
\usepackage{mathtools}
\usepackage{amsthm}
\usepackage{thmtools}
\usepackage[unq]{unique}
\DeclareMathOperator{\powerset}{\mathcal{P}}

\usepackage[binary-units]{siunitx}

\usepackage{pgf}

\usepackage[colorlinks]{hyperref}
\usepackage[capitalize]{cleveref}
\crefalias{question}{item}
\crefname{question}{question}{questions}
\Crefname{question}{Question}{Questions}
\creflabelformat{question}{#2\textup{#1}#3}

% adapted from https://tex.stackexchange.com/a/30726/17418
\newcommand\chnote[1]{%
  \begingroup
  \renewcommand\thefootnote{}\footnote{#1}%
  \addtocounter{footnote}{-1}%
  \endgroup
}

\usepackage[marginparmargin=right]{didactic}
 after
% \documentclass but before \begin{document}.

%\usepackage[utf8]{inputenc}
\usepackage[T1]{fontenc}
\usepackage[swedish,british]{babel}
\usepackage{booktabs}

\usepackage{noweb}
% Caveat about longxref: it typesets "This definition is continued in
% chunks ..." lists whose line count depends on the page numbers of later
% chunks.  With many continuations (e.g. <<test functions>>) the list
% length can feed back into the pagination it describes, so the labels
% never converge and latexmk reruns forever ("Label(s) may have
% changed.").  If that happens, adjust the prose near the oscillating
% page boundary (diff two consecutive .aux files to find it) until the
% layout has a fixed point again.
\noweboptions{shift,breakcode,longxref,longchunks}

\usepackage[natbib,style=alphabetic,maxbibnames=99]{biblatex}
\addbibresource{bibliography.bib}

\usepackage[inline]{enumitem}
\setlist[enumerate]{label=(\arabic*)}

\usepackage[all]{foreign}
\renewcommand{\foreignfullfont}{}
\renewcommand{\foreignabbrfont}{}

%\usepackage{newclude}
\usepackage{import}

\usepackage[strict]{csquotes}
\usepackage[single]{acro}

\usepackage{subcaption}
\usepackage{graphicx}

\usepackage[noend]{algpseudocode}
\usepackage{xparse}

\let\email\texttt

%\usepackage[outputdir=ltxobj]{minted}
\usepackage{minted}
\setminted{autogobble,linenos}

\usepackage{pythontex}
\setpythontexoutputdir{.}
\setpythontexworkingdir{..}

\usepackage{fancyvrb}

\usepackage{amsmath}
\usepackage{amssymb}
\usepackage{mathtools}
\usepackage{amsthm}
\usepackage{thmtools}
\usepackage[unq]{unique}
\DeclareMathOperator{\powerset}{\mathcal{P}}

\usepackage[binary-units]{siunitx}

\usepackage{pgf}

\usepackage[colorlinks]{hyperref}
\usepackage[capitalize]{cleveref}
\crefalias{question}{item}
\crefname{question}{question}{questions}
\Crefname{question}{Question}{Questions}
\creflabelformat{question}{#2\textup{#1}#3}

% adapted from https://tex.stackexchange.com/a/30726/17418
\newcommand\chnote[1]{%
  \begingroup
  \renewcommand\thefootnote{}\footnote{#1}%
  \addtocounter{footnote}{-1}%
  \endgroup
}

\usepackage[marginparmargin=right]{didactic}
 after
% \documentclass but before \begin{document}.

%\usepackage[utf8]{inputenc}
\usepackage[T1]{fontenc}
\usepackage[swedish,british]{babel}
\usepackage{booktabs}

\usepackage{noweb}
% Caveat about longxref: it typesets "This definition is continued in
% chunks ..." lists whose line count depends on the page numbers of later
% chunks.  With many continuations (e.g. <<test functions>>) the list
% length can feed back into the pagination it describes, so the labels
% never converge and latexmk reruns forever ("Label(s) may have
% changed.").  If that happens, adjust the prose near the oscillating
% page boundary (diff two consecutive .aux files to find it) until the
% layout has a fixed point again.
\noweboptions{shift,breakcode,longxref,longchunks}

\usepackage[natbib,style=alphabetic,maxbibnames=99]{biblatex}
\addbibresource{bibliography.bib}

\usepackage[inline]{enumitem}
\setlist[enumerate]{label=(\arabic*)}

\usepackage[all]{foreign}
\renewcommand{\foreignfullfont}{}
\renewcommand{\foreignabbrfont}{}

%\usepackage{newclude}
\usepackage{import}

\usepackage[strict]{csquotes}
\usepackage[single]{acro}

\usepackage{subcaption}
\usepackage{graphicx}

\usepackage[noend]{algpseudocode}
\usepackage{xparse}

\let\email\texttt

%\usepackage[outputdir=ltxobj]{minted}
\usepackage{minted}
\setminted{autogobble,linenos}

\usepackage{pythontex}
\setpythontexoutputdir{.}
\setpythontexworkingdir{..}

\usepackage{fancyvrb}

\usepackage{amsmath}
\usepackage{amssymb}
\usepackage{mathtools}
\usepackage{amsthm}
\usepackage{thmtools}
\usepackage[unq]{unique}
\DeclareMathOperator{\powerset}{\mathcal{P}}

\usepackage[binary-units]{siunitx}

\usepackage{pgf}

\usepackage[colorlinks]{hyperref}
\usepackage[capitalize]{cleveref}
\crefalias{question}{item}
\crefname{question}{question}{questions}
\Crefname{question}{Question}{Questions}
\creflabelformat{question}{#2\textup{#1}#3}

% adapted from https://tex.stackexchange.com/a/30726/17418
\newcommand\chnote[1]{%
  \begingroup
  \renewcommand\thefootnote{}\footnote{#1}%
  \addtocounter{footnote}{-1}%
  \endgroup
}

\usepackage[marginparmargin=right]{didactic}
 after
% \documentclass but before \begin{document}.

%\usepackage[utf8]{inputenc}
\usepackage[T1]{fontenc}
\usepackage[swedish,british]{babel}
\usepackage{booktabs}

\usepackage{noweb}
% Caveat about longxref: it typesets "This definition is continued in
% chunks ..." lists whose line count depends on the page numbers of later
% chunks.  With many continuations (e.g. <<test functions>>) the list
% length can feed back into the pagination it describes, so the labels
% never converge and latexmk reruns forever ("Label(s) may have
% changed.").  If that happens, adjust the prose near the oscillating
% page boundary (diff two consecutive .aux files to find it) until the
% layout has a fixed point again.
\noweboptions{shift,breakcode,longxref,longchunks}

\usepackage[natbib,style=alphabetic,maxbibnames=99]{biblatex}
\addbibresource{bibliography.bib}

\usepackage[inline]{enumitem}
\setlist[enumerate]{label=(\arabic*)}

\usepackage[all]{foreign}
\renewcommand{\foreignfullfont}{}
\renewcommand{\foreignabbrfont}{}

%\usepackage{newclude}
\usepackage{import}

\usepackage[strict]{csquotes}
\usepackage[single]{acro}

\usepackage{subcaption}
\usepackage{graphicx}

\usepackage[noend]{algpseudocode}
\usepackage{xparse}

\let\email\texttt

%\usepackage[outputdir=ltxobj]{minted}
\usepackage{minted}
\setminted{autogobble,linenos}

\usepackage{pythontex}
\setpythontexoutputdir{.}
\setpythontexworkingdir{..}

\usepackage{fancyvrb}

\usepackage{amsmath}
\usepackage{amssymb}
\usepackage{mathtools}
\usepackage{amsthm}
\usepackage{thmtools}
\usepackage[unq]{unique}
\DeclareMathOperator{\powerset}{\mathcal{P}}

\usepackage[binary-units]{siunitx}

\usepackage{pgf}

\usepackage[colorlinks]{hyperref}
\usepackage[capitalize]{cleveref}
\crefalias{question}{item}
\crefname{question}{question}{questions}
\Crefname{question}{Question}{Questions}
\creflabelformat{question}{#2\textup{#1}#3}

% adapted from https://tex.stackexchange.com/a/30726/17418
\newcommand\chnote[1]{%
  \begingroup
  \renewcommand\thefootnote{}\footnote{#1}%
  \addtocounter{footnote}{-1}%
  \endgroup
}

\usepackage[marginparmargin=right]{didactic}
